\documentclass[tikz,border=10pt]{standalone}
\usepackage{tikz}
\usepackage{amsmath}
\usepackage{amssymb}
\usepackage{ctex}
\usepackage{pgfplots}
\pgfplotsset{compat=1.18}
\usetikzlibrary{calc, positioning, arrows.meta, decorations.pathreplacing}

\begin{document}
\begin{tikzpicture}[>=Stealth]

%% ===== 条形图 + 累积能量曲线 =====
\begin{axis}[
  name=axbar,
  width=13cm, height=8cm,
  xmin=0.3, xmax=10.7,
  ymin=0, ymax=115,
  axis lines=left,
  xlabel={奇异值 index $i$},
  ylabel={$\sigma_i$},
  xlabel style={font=\small},
  ylabel style={font=\small},
  xtick={1,2,3,4,5,6,7,8,9,10},
  xticklabels={$1$,$2$,$3$,$4$,$5$,$6$,$7$,$8$,$9$,$n$},
  ytick={0,20,40,60,80,100},
  yticklabels={$0$,$20$,$40$,$60$,$80$,$100\%$},
  tick label style={font=\scriptsize},
  grid=both, grid style={gray!12},
  clip=false,
  % 第二y轴（累积能量%）用同一坐标范围 0-100 映射到 0-100
]

  %% 奇异值条形（指数衰减: sigma_i = 90 * 0.55^(i-1)）
  % i=1: 90.0, i=2: 49.5, i=3: 27.2, i=4: 15.0, i=5: 8.2
  % i=6: 4.5, i=7: 2.5, i=8: 1.4, i=9: 0.75, i=10: 0.41

  % 前 k=5 个（蓝色突出）
  \addplot[ybar, bar width=0.55cm, fill=blue!70!black, draw=blue!40,
           bar shift=0]
    coordinates {(1,90.0) (2,49.5) (3,27.2) (4,15.0) (5,8.2)};

  % 后 n-k=5 个（浅色）
  \addplot[ybar, bar width=0.55cm, fill=gray!30, draw=gray!40,
           bar shift=0]
    coordinates {(6,4.5) (7,2.5) (8,1.4) (9,0.75) (10,0.41)};

  %% 分界虚线 (k=5 和 k+1=6 之间)
  \addplot[dashed, gray!60, thick, domain=5.5:5.5]
    coordinates {(5.5,0) (5.5,100)};

  %% 累积能量曲线（虚线，右轴%）
  % total = 90+49.5+27.2+15+8.2+4.5+2.5+1.4+0.75+0.41 = 199.46
  % cumsum: 90/199.46=45.1%, 139.5/199.46=69.9%, 166.7/199.46=83.6%
  %          181.7/199.46=91.1%, 189.9/199.46=95.2%
  %          194.4/199.46=97.5%, 196.9/199.46=98.7%, 198.3/199.46=99.4%
  %          199.05/199.46=99.8%, 199.46/199.46=100%
  \addplot[red!70!black, dashed, very thick, mark=o, mark size=2.5pt,
           mark options={fill=red!60}]
    coordinates {
      (1,45.1) (2,69.9) (3,83.6) (4,91.1) (5,95.2)
      (6,97.5) (7,98.7) (8,99.4) (9,99.8) (10,100.0)
    };

  %% k=5 分界标注
  \node[font=\scriptsize, blue!80!black, rotate=90]
    at (axis cs:5.5, 55) {\ $k=5$};

  %% 前k标注
  \node[font=\scriptsize, fill=blue!10, draw=blue!40, rounded corners=2pt,
        inner sep=3pt, align=center]
    at (axis cs:2.5, 75) {前 $k$ 个\\蓝色保留};

  %% 后n-k标注
  \node[font=\scriptsize, fill=gray!15, draw=gray!40, rounded corners=2pt,
        inner sep=3pt, align=center]
    at (axis cs:8.0, 40) {后 $n-k$ 个\\截断丢弃};

  %% 累积曲线标注
  \node[font=\scriptsize, red!80!black, fill=red!8, draw=red!40,
        rounded corners=2pt, inner sep=3pt]
    at (axis cs:7.5, 80) {累积能量曲线};

  %% 95%线
  \addplot[orange!70, dashed, thin, domain=0.5:5.3]
    {95.2};
  \node[font=\scriptsize, orange!80!black, right]
    at (axis cs:5.4, 95.2) {$95.2\%$（前$5$个）};

\end{axis}

%% ===== 图例（轴外）=====
\node[font=\small, fill=white, draw=gray!50, thin, rounded corners=4pt,
      inner sep=8pt, align=left]
  at ([xshift=0cm, yshift=-2.8cm]axbar.south) {
    \rule{12pt}{8pt}\textcolor{blue!70!black}{\rule{12pt}{8pt}}\ 保留的前 $k$ 个奇异值
    \quad
    \textcolor{gray!50}{\rule{12pt}{8pt}}\ 截断的后 $n-k$ 个
    \quad
    \textcolor{red!70!black}{- - -}\ 累积能量 $\displaystyle E_k=\frac{\sum_{i=1}^k\sigma_i^2}{\sum_{i=1}^n\sigma_i^2}$
  };

%% ===== 底部公式框 =====
\node[font=\normalsize, fill=white, draw=gray!50, thin, rounded corners=5pt,
      inner sep=10pt, align=center]
  at ([xshift=0cm, yshift=-4.8cm]axbar.south) {
    \textbf{截断 SVD（Truncated SVD）：}
    $\displaystyle A \approx \sum_{i=1}^{k} \sigma_i\, \mathbf{u}_i\mathbf{v}_i^T
    = U_k \Sigma_k V_k^T$
    \quad（秩-$k$ 最佳逼近，Eckart–Young 定理）
  };

%% ===== 整体标题 =====
\node[font=\large\bfseries] at ([yshift=0.8cm]axbar.north)
  {截断 SVD：奇异值能量分布与保留比例};

\end{tikzpicture}
\end{document}
